% Section VII (plan v2): certificates that bound events. Drafted 2026-09-06; subsections A-C are added as their results arrive (feat-019/020/021); D uses results/warped_anchor.csv (feat-022).
\section{What Would Make the Certificate Say More}\label{sec:certificates}

The gaps found so far are properties of the guarantee, not of the implementation, and each admits a repair inside the decoder. This section implements three and prices them: a pathwise budget on the realised log-ratio, which certifies every event (Section~\ref{sec:pathwise}), priced against the unconstrained model by a judge outside the audit (Section~\ref{sec:utility}); a concentration certificate for the unchanged KL decoder, computed from quantities the solver already evaluates (Section~\ref{sec:concentration}); and a per-user budget with a cap on the bank (Section~\ref{sec:odometer}). Sections~\ref{sec:warped} and~\ref{sec:leakage} then show what all three, and Theorem~3.1 itself, are stated relative to: an anchor that the deployer's decoding settings reshape, and that has already read the most quoted lines of the works it is meant to protect.

\subsection{A pathwise budget certifies events}\label{sec:pathwise}
The KL certificate bounds an expectation because the decoder charges one: $a_t = D_{\mathrm{KL}}(p_\theta \,\|\, p_s)$ averages the log-ratio over tokens that were not sampled. Charging instead the log-ratio of the sampled token, $r_t = \log p_\theta(y_t) - \log p_s(y_t)$, gives a decoder whose invariant is the realised ratio itself. At every step it takes the largest $\theta_t \in [0, 1]$ whose maximal log-ratio $m_t(\theta) = \max_v \log p_\theta(v)/p_s(v)$ fits the allowance $k^{\mathrm{pw}}_t = \max\{0, (t+1)k - \delta_{\mathrm{init}}(x) - R_{t-1}\}$, with $R_{t-1} = \sum_{i < t} r_i$; $m_t$ increases in $\theta$ from $m_t(0) = 0$, so a bisection finds $\theta_t$, and $\theta_t = 1$ whenever the risky model fits. Proposition~\ref{prop:pathwise} shows that every output then satisfies $L(y) \le K - \delta_{\mathrm{init}}(x)$, the $\Delta_{\max}$ form of near-access-freeness, so $p_\theta(E \mid x) \le e^{K} p_s(E \mid x)$ for every event $E$ by Lemma~2.2 of Vyas et al., with the KL form as a corollary. The banking rule is unchanged, a token the anchor preferred refunds allowance, and the per-path sum of KL spends is no longer bounded and need not be (it reaches $1.12K$ at $k=1$ while its expectation stays below $K$).

We ran the composition attack against the pathwise decoder on the same passages, seeds, and budgets. The invariant held on all $23{,}844$ budgeted queries against the memorising 8B model and on the $1{,}250$ queries of the 70B audit at the authors' settings, with no overshoot. The event bound costs the adversary a little at every budget: with the memorising model, oracle windows recover $0.06$, $0.19$, $0.39$, and $0.73$ of a passage at $k = 3$, $5$, $10$, and $20$ against $0.10$, $0.23$, $0.41$, and $0.79$ under the KL decoder, chained windows $0.00$, $0.02$, $0.09$, and $0.44$ against $0.01$, $0.03$, $0.11$, and $0.50$, and a single query $0.45$ against $0.48$ at $k=20$; at $k=50$ both return the unconstrained model. The difference is small because the decoders differ only on the few tokens whose maximal log-ratio exceeds the average by more than the allowance absorbs. What changes is what can be said: the same numbers are now guaranteed for every output, and the reconstruction rate of any window is bounded by $\exp(K_i - S_i)$ (Proposition~\ref{prop:cost}). With the 70B model at the authors' settings the pathwise decoder holds a single query to $0.02$, $0.05$, $0.09$, and $0.31$ and oracle windows to $0.02$, $0.07$, $0.22$, and $0.55$ at $k = 3$, $5$, $10$, and $20$, against $0.31$ and $0.61$ for the sampled model alone.

The price is paid on ordinary prompts. On the sweep of Section~\ref{sec:sweep}, with both decoders on the same $900$ prompts, seeds, and budgets ($2{,}700$ trajectories per budget), the pathwise decoder intervenes more often at every budget: the constraint is active on $56\%$ of decode steps against $43\%$ at $k=0.5$, $14.6\%$ against $5.0\%$ at $k=1$, $1.9\%$ against $0.3\%$ at $k=3$, $1.1\%$ against $0.2\%$ at $k=5$, and $0.5\%$ against $0.01\%$ at $k=10$, and it forces the anchor on $6.1\%$ of steps against $2.2\%$ at $k=3$ (at $k \le 1$ the forced fraction is the prefix debt's and the same for both). The risky model is served unchanged on $92.0\%$ of steps against $97.6\%$ at $k=3$ and on $99.5\%$ against $99.9\%$ at $k=10$; at $k=20$ the two decoders coincide. All $16{,}200$ pathwise trajectories satisfied $L(y) \le \max\{0, B\}$. What those extra interventions cost is not something the risky model can be asked, since it is the distribution that moved; Section~\ref{sec:utility} measures it from outside.

Proposition~\ref{prop:cost} bounds what re-sampling buys the adversary: a window with anchor surprisal $S_i$ and budget $K_i = kL$ is reproduced exactly with probability at most $e^{K_i - S_i}$ per attempt. Measured against an adversary who re-samples each $50$-token window up to eight times, the bound held wherever it was informative, and no window was reproduced exactly at $k \le 3$ in $2{,}376$ attempts (Appendix~\ref{app:tables}). It is vacuous at the budgets where the attack succeeds, because a $50$-token window costs the anchor about $150$ nats and $K_i$ is already $250$ at $k=5$: the cost bound and the certificate of Section~\ref{sec:certificate} go vacuous at the same budget, for the same reason, and below that budget both bind. Where the bound is empty the adversary's cost is set by the budget binding instead, at $3.8$, $38$ and $153$ exact reproductions per thousand attempts at $k=5$, $10$ and $20$.

\subsection{What the constraint costs}\label{sec:utility}
A decoder scored by the log-loss of the model whose distribution it altered cannot report its own price, so we measured it externally. The sweep of Section~\ref{sec:sweep} already contains both decoders and both baselines on the same $500$ ordinary prompts and seeds across the general-knowledge, biography and creative classes, $1{,}500$ trajectories per budget. An instruction-tuned model outside the audit, Qwen2.5-7B-Instruct, judged $600$ completions per budget against the unconstrained risky model's on the same prompt, with the presentation order randomised per pair. The number that makes the rest readable is the null: the unconstrained model judged against \emph{itself} at a different seed loses $47.0\%$ of the time, because temperature-1 sampling and an imperfect judge already disagree. Every row of Table~\ref{tab:utility} is a deviation from that, not from a half.

\begin{table}[t]
\centering
\caption{The price of the constraint on $1{,}500$ ordinary generations per budget. ``Active'' is the fraction of decode steps the constraint touched and ``loss'' the fraction of $600$ judged pairs the arm lost to the unconstrained model on the same prompt. The unconstrained model against itself loses $47.0\%$; the anchor alone loses $62.7\%$. Distinct-3 per arm is in the artifact. Differences under about $4$ points are within two standard errors.}
\label{tab:utility}
\small
\setlength{\tabcolsep}{3.4pt}
\begin{tabular}{@{}lrrrr@{}}
\toprule
& \multicolumn{2}{c}{KL decoder} & \multicolumn{2}{c}{pathwise decoder} \\
\cmidrule(lr){2-3}\cmidrule(l){4-5}
$k$ & active & loss & active & loss \\
\midrule
0.5 & 48.0\% & 59.8\% & 63.3\% & 60.7\% \\
1 & 5.43\% & 52.3\% & 16.1\% & 50.8\% \\
3 & 0.35\% & 47.5\% & 2.08\% & 51.7\% \\
5 & 0.24\% & 45.7\% & 1.15\% & 49.8\% \\
10 & 0.01\% & 43.0\% & 0.51\% & 45.7\% \\
20 & 0.00\% & 44.8\% & 0.00\% & 41.5\% \\
\bottomrule
\end{tabular}
\end{table}

Two readings follow, and they are the same fact seen twice. At $k \ge 3$ neither decoder is distinguishable from not constraining at all: the KL decoder loses $47.5$, $45.7$, $43.0$ and $44.8$ per cent at $k = 3, 5, 10, 20$ against the null's $47.0$, because at those budgets it intervenes on at most a third of one per cent of steps. The budgets the mechanism's authors evaluate are free because they are inert. The cost appears only where Section~\ref{sec:certificate} found the certificate informative: at $k=1$ the KL decoder loses $52.3\%$ and at $k=0.5$ it loses $59.8\%$, against $62.7\%$ for serving the anchor alone, so a budget small enough to bound reproduction has already given up most of what separates the risky model from a 1.8B model that never read the book. The reference-free measure agrees without a judge, distinct-3 falling from $0.985$ to $0.856$ over the same range. The event-level guarantee is not free either: in the deployed band the pathwise decoder loses more than the KL decoder at the same budget, $51.7\%$ against $47.5\%$ at $k=3$ and $49.8\%$ against $45.7\%$ at $k=5$, each about two standard errors, converging above $k=10$ where both are inert. That is the price of Proposition~\ref{prop:pathwise}: small in absolute terms, and paid exactly where the KL decoder pays nothing.

\subsection{What the KL certificate concentrates to}\label{sec:concentration}
A deployer who keeps the KL decoder can still say something about single outputs, because the gap between the realised ratio and the charged spend is a martingale whose increments the decoder already computes. Write $M_t = \sum_{i \le t}(r_i - a_i)$; then $L(y) = Z_T + M_T$ with $Z_T \le K$ on every path, so an output exceeds the certificate by $\lambda$ only if the martingale does. Each increment is bounded above by $b_i = m_i - a_i$ and has conditional variance $\sigma_i^2 = \mathrm{Var}_{p_{\theta_i}}[r_i]$, both functions of the prefix evaluated before the token is drawn, so Proposition~\ref{prop:conc} turns Freedman's inequality~\cite{freedman1975tail,tropp2011freedman} into a certificate the deployer can enforce by clipping them, or verify after the fact from the logs. This is the concentration assumption that Lemma~2.3 of Vyas et al.\ requires, stated in terms the decoder controls. Evaluated on the sweep at the 90th percentiles of the logged $\sum_i \sigma_i^2$ and $\max_i b_i$, it gives $\Pr[L(y) \ge 2K] \le 0.004$ at $k=0.5$ and $\Pr[L(y) \ge K+100] \le 0.011$ at every $k \ge 3$, against empirical frequencies of zero (Appendix~\ref{app:theory} tabulates the inputs). The certificate is therefore informative and loose by two orders of magnitude, because Freedman's inequality pays on every path for a worst-case increment of $15$ to $30$ nats that few tokens produce. What it cannot do is repair the budgets where the problem lies: $\delta(0) = 1$, so it says nothing about the $34\%$ ($k=0.5$) and $8\%$ ($k=1$) of outputs that exceed $K$ itself. Concentration does not rescue the KL form at the budgets where that form is informative; bounding $b_i$ per step does, which is what the pathwise decoder enforces.

\subsection{Accounting across queries}\label{sec:odometer}
Near-access-freeness does not compose~\cite{cohen2025blameless}, and Section~\ref{sec:attack-results} shows the consequence. Differential privacy meets the same problem with a per-user budget tracked by an odometer~\cite{rogers2016odometers,whitehouse2023fully}, and the device transfers because the spend of every query is a number the decoder already computes. We replayed the logs of the composition attack (memorising model, 100 passages) against a per-user budget $B_{\mathrm{user}}$: queries are served in the attack's order until the cumulative spend would exceed the budget, the user is then refused, and reconstruction is scored on the outputs served (Table~\ref{tab:odometer}).

\begin{figure}[t]
\centering
\includegraphics[width=\columnwidth]{separation.pdf}
\caption{What a per-user budget buys. Horizontal axis: the number of ordinary $200$-token completions the same budget admits to a reader, at the measured $0.82$ nats per token. Vertical axis: the fraction of a protected passage the strongest strategy at that $k$ recovers before the budget is exhausted. At $k=20$ the curve has crossed a tenth of the work before the budget admits one ordinary answer; at $k=3$ it never crosses.}
\label{fig:separation}
\end{figure}

The odometer works: reconstructing a passage costs the attack $586$ nats at $k=3$ and $860$ at $k=20$ against a median anchor surprisal of $849$, so a budget of $400$ nats bounds reconstruction to $30\%$ (oracle) and $26\%$ (chained) and $200$ nats bounds it to $12\%$. The question a deployer has to answer is whether any budget is at once permissive to a reader and restrictive to a reconstructor, and Proposition~\ref{prop:sep} says it is not, for a reason that is not a matter of tuning. Both activities are metered in the same additive currency, so a budget $B$ admits about $B/c_{\mathrm{legit}}$ tokens of ordinary completion and bounds recovery at about $B/c_{\mathrm{rec}}(k)$ tokens of the work, and the filter's whole protective power is the ratio $c_{\mathrm{rec}}(k)/c_{\mathrm{legit}}$, which $k$ fixes. We measured both rates on the same decoder. Ordinary traffic costs $c_{\mathrm{legit}} = 0.82$ nats per token, and remarkably evenly: the per-class medians across all six prompt classes and every budget lie between $0.73$ and $0.90$, so one $200$-token answer costs $164$ nats. Reconstruction is dearer per token recovered, and the ratio collapses as the decoder approaches the identity --- $13.6$ at $k=5$, $8.6$ at $k=10$, $4.56$ at $k=20$.

That collapse is the problem, because the budget is redundant wherever it is comfortable and uncomfortable wherever it binds (Fig.~\ref{fig:separation}). At $k \le 3$ no strategy recovers a tenth of a work even with no budget at all, so the filter has nothing to do. At $k=20$, the top of the range the mechanism's authors sweep and the only part of it where the decoder leaves the risky model's utility intact, holding the strongest adversary to a tenth of a work requires $B = 100$ nats: $0.61$ of a single ordinary completion. The proposition is stated for filters that meter the decoder's own additive spend, and Remark~\ref{rem:sep} says what escapes it: a semantic detector, a per-work rather than per-user budget, or any accounting that charges the two activities at different rates.

\paragraph{Capping the bank}
The banking rule of Eq.~\eqref{eq:banking} is a token bucket of unbounded depth, and Section~\ref{sec:attack-results} showed that at $k=10$ the bank, rather than the rate, pays for most of a memorised passage. A cap $D$ on the bank discards allowance above $D$, so no $W$ consecutive steps can spend more than $Wk + D$, and $D = k$ is a plain rate limit. We decoded the attack again with $D = k$ and $D = 5k$ on the same seeds ($3{,}804$ queries, all within budget). At $k=10$ the rate limit cuts single-query reconstruction from $0.20$ to $0.12$ and chained windows from $0.11$ to $0.07$, $D = 5k$ gives $0.17$ and $0.09$, and oracle windows stay at $0.39$ to $0.41$ at every depth, because a $50$-token window at $10$ nats per step pays for its tokens without saving. At $k=20$ the cap changes nothing ($0.46$, $0.80$, and $0.47$ against $0.48$, $0.79$, and $0.50$), as the budget-path analysis predicts, since the rate alone exceeds the anchor's surprisal on $92\%$ of passages. The cap bounds bursts, which a deployer may want for its own sake, but not reconstruction.




\subsection{The certificate depends on the decoding settings}\label{sec:warped}
He et al.\ apply the temperature and the repetition penalty to both models' logits before the KL solve (their Appendix~B) and decode books at temperature $0.7$ with penalty $1.1$ (their Appendix~D.1), so Theorem~3.1 certifies proximity to the \emph{warped} anchor $p_s^{(\tau,\rho)}$ and the surprisal in Eq.~\eqref{eq:cap} must be computed under it (Remark~\ref{rem:warp}). Tempering sharpens the anchor, so the median surprisal of a CopyBench passage rises from $3.20$ nats per token at temperature $1$ to $3.55$ at the authors' setting and $6.53$ at temperature $0.3$, and the vacuous fraction at $k=1$ falls from $44\%$ to $23\%$ and then to $0\%$; at $k \ge 3$ the certificate stays vacuous for every passage under every setting but the last (Appendix~\ref{app:tables}). A published $K$ therefore means a different bound under settings the deployer chooses after the theorem is proved, which is why they belong on the certificate.

\subsection{What the anchor already knows}\label{sec:leakage}
Near-access-freeness is relative to the anchor: text the anchor assigns high probability to is protected by no $K$, and He et al.\ note that famous fragments reach openly licensed corpora through quotation (their Appendix~A.3). We measured the leak. Across 50 public-domain novels the anchor's per-token surprisal of the opening 100 tokens has median $3.11$ nats against $3.33$ for a passage $20{,}000$ characters into the same book, and the two most quoted openings in English cost $1.2$ nats per token; for the sixteen CopyBench novels the widely quoted first sentence costs $4.35$ nats per token against $4.73$ for a passage of the same length from the same book ($2.8$ for \emph{Harry Potter}). The leak stops short of reproduction, since the anchor alone recovers none of the 50 continuations, and Eq.~\eqref{eq:cap} charges it correctly by lowering $S(x)$ only where the text is famous. It is nevertheless a fact about the anchor's corpus rather than about the mechanism, which is the reason a certificate card should report the anchor's surprisal on the reference set.
